% Bagian: first-order-logic
% Bab: beyond
% Subbagian: intuitionistic-logic

\documentclass[../../../include/open-logic-section]{subfiles}

\begin{document}

\olfileid[id]{fol}{byd}{il}

\olsection{Logika Intuisionistik}

Berbeda dari logika orde kedua dan logika orde tinggi, logika orde pertama
intuisionistik merupakan pembatasan atas versi klasik, yang dimaksudkan untuk
memodelkan jenis penalaran yang lebih ``konstruktif''. Contoh-contoh berikut
dapat menggambarkan beberapa motivasi yang mendasarinya.

Andaikan suatu hari seseorang menghampiri Anda dan mengumumkan bahwa ia telah
menentukan suatu bilangan asli~$x$ dengan sifat bahwa jika $x$ prima, hipotesis
Riemann benar, dan jika $x$ komposit, hipotesis Riemann salah. Kabar bagus!{}
Benar atau tidaknya hipotesis Riemann merupakan salah satu pertanyaan terbuka
besar dalam matematika, dan di sini orang itu tampaknya telah mereduksi
masalahnya menjadi masalah perhitungan, yakni menentukan apakah suatu bilangan
tertentu prima atau bukan.

Berapakah nilai ajaib~$x$? Orang itu menguraikannya sebagai berikut: $x$
adalah bilangan asli yang sama dengan~$7$ jika hipotesis Riemann benar, dan
sama dengan~$9$ jika tidak.

Dengan marah, Anda menuntut uang Anda kembali. Dari sudut pandang klasik,
uraian di atas memang menentukan satu nilai~$x$ yang unik; tetapi yang
sesungguhnya Anda inginkan ialah nilai~$x$ yang diberikan secara
\emph{eksplisit}.

Sebagai contoh lain yang mungkin tidak terlalu dibuat-buat, pertimbangkan
pertanyaan berikut. Kita tahu bahwa suatu bilangan irasional dapat dipangkatkan
dengan bilangan rasional dan menghasilkan bilangan rasional. Misalnya,
$\sqrt{2}^2 = 2$. Yang kurang jelas ialah apakah suatu bilangan irasional dapat
dipangkatkan dengan bilangan \emph{irasional} dan menghasilkan bilangan
rasional. Teorema berikut memberikan jawaban afirmatif:

\begin{thm}
Terdapat bilangan irasional $a$ dan $b$ sedemikian sehingga $a^b$ rasional.
\end{thm}

\begin{proof}
Pertimbangkan $\sqrt{2}^{\sqrt{2}}$. Jika bilangan ini rasional, kita selesai:
kita dapat mengambil $a = b = \sqrt{2}$. Jika tidak, bilangan itu irasional.
Selanjutnya,
\[
(\sqrt{2}^{\sqrt{2}})^{\sqrt{2}} = \sqrt{2}^{\sqrt{2} \cdot
  \sqrt{2}} = \sqrt{2}^2 = 2,
\]
yang tentu saja rasional. Jadi, dalam kasus ini, ambil $a$ sebagai
$\sqrt{2}^{\sqrt{2}}$, dan ambil $b$ sebagai~$\sqrt 2$.
\end{proof}

Apakah ini merupakan bukti yang valid? Kebanyakan matematikawan merasa
demikian. Namun, sekali lagi, ada sesuatu yang agak tidak memuaskan di sini:
kita telah membuktikan keberadaan sepasang bilangan real dengan sifat tertentu
tanpa dapat mengatakan pasangan bilangan \emph{mana} yang dimaksud. Hasil yang
sama dapat dibuktikan dengan cara yang \emph{memberikan} pasangan $a$, $b$ di dalam
bukti itu sendiri: ambil $a = \sqrt{3}$ dan $b = \log_3 4$. Maka
\[
a^b = \sqrt{3}^{\log_3 4} = 3^{1/2 \cdot \log_3 4} = (3^{\log_3
  4})^{1/2} = 4^{1/2}= 2,
\]
karena $3^{\log_3 x} = x$.

Logika intuisionistik dirancang untuk memodelkan jenis penalaran yang tidak
memperbolehkan langkah seperti dalam bukti pertama. Membuktikan keberadaan
suatu~$x$ yang memenuhi~$!A(x)$ berarti kita harus memberikan~$x$ tertentu
dan bukti bahwa nilai itu memenuhi~$!A$, seperti dalam bukti kedua. Membuktikan
bahwa $!A$ atau $!B$ berlaku mengharuskan kita dapat membuktikan salah satunya.

Secara formal, logika orde pertama intuisionistik diperoleh dengan membatasi
sistem !!a{derivation} bagi logika orde pertama dengan cara tertentu. Demikian
pula, terdapat versi intuisionistik dari logika orde kedua maupun orde tinggi.
Dari sudut pandang matematis, semua itu hanyalah sistem deduktif formal,
tetapi, sebagaimana telah dicatat, sistem-sistem tersebut dimaksudkan untuk
memodelkan suatu jenis penalaran matematis. Orang dapat menganggapnya sebagai
jenis penalaran yang dibenarkan oleh pandangan filosofis tertentu tentang
matematika (seperti intuisionisme Brouwer); orang dapat menganggapnya sebagai
jenis penalaran matematis yang lebih ``konkret'' dan memuaskan (sejalan dengan
konstruktivisme Bishop); dan orang dapat memperdebatkan apakah uraian formal
tersebut menangkap motivasi informalnya. Namun, apa pun pendirian filosofis
kita, kita dapat mengkaji logika intuisionistik sebagai logika yang disajikan
secara formal; dan apa pun alasannya, banyak ahli logika matematika tertarik
untuk melakukannya.

Terdapat suatu interpretasi konstruktif informal bagi penghubung
intuisionistik, yang biasanya dikenal sebagai interpretasi BHK (dinamai
menurut Brouwer, Heyting, dan Kolmogorov). Interpretasi itu berbunyi sebagai
berikut: suatu bukti dari $!A \land !B$ terdiri atas bukti dari $!A$ yang
dipasangkan dengan bukti dari $!B$; suatu bukti dari $!A \lor !B$ terdiri atas
bukti dari $!A$ atau bukti dari $!B$, dengan informasi eksplisit tentang mana
yang berlaku; suatu bukti dari $!A \lif !B$ terdiri atas prosedur yang
mengubah bukti dari $!A$ menjadi bukti dari~$!B$; suatu bukti dari
$\lforall[x][!A(x)]$ terdiri atas prosedur yang mengembalikan bukti dari
$!A(x)$ untuk setiap nilai~$x$; dan suatu bukti dari $\lexists[x][!A(x)]$
terdiri atas suatu nilai~$x$, beserta bukti bahwa nilai itu memenuhi~$!A$.
Kita dapat menguraikan interpretasi tersebut dalam istilah komputasional yang
dikenal sebagai ``isomorfisme Curry--Howard'' atau ``paradigma
!!{formula}s-sebagai-tipe'': pandanglah !!a{formula} sebagai penentu suatu
jenis tipe data, dan bukti sebagai objek komputasional dari tipe data tersebut
yang memungkinkan kita melihat bahwa !!{formula} yang bersesuaian benar.

Logika intuisionistik sering dipandang sebagai logika klasik ``dikurangi''
hukum tiada jalan tengah. Teorema berikut membuat pernyataan ini lebih tepat.
\begin{thm}
Secara intuisionistik, skema-skema aksioma berikut ekuivalen:
\begin{enumerate}
\item $(\lnot !A \lif \lfalse) \lif !A$.
\item $!A \lor \lnot !A$
\item $\lnot \lnot !A \lif !A$
\end{enumerate}
\end{thm}
Memperoleh contoh setiap skema dari salah satu di antara dua skema lainnya
merupakan latihan yang baik dalam logika intuisionistik.

Sistem deduktif pertama bagi logika proposisional intuisionistik, yang
diajukan sebagai formalisasi intuisionisme Brouwer, dikemukakan secara
terpisah oleh Kolmogorov, Glivenko, dan Heyting. Formalisasi pertama logika
orde pertama intuisionistik (serta sebagian matematika intuisionistik)
dikemukakan oleh Heyting. Meskipun sejumlah skema yang valid secara klasik
tidak valid secara intuisionistik, banyak pula yang valid.

\emph{Terjemahan negasi ganda} menguraikan suatu hubungan penting antara
logika klasik dan logika intuisionistik. Terjemahan itu didefinisikan secara
induktif sebagai berikut (pandanglah $!A^N$ sebagai terjemahan
``intuisionistik'' dari !!{formula} klasik~$!A$):
\begin{align*}
!A^N & \ident \lnot\lnot !A \quad \text{untuk !!{formula}s atomik $!A$} \\
(!A \land !B)^N & \ident (!A^N \land !B^N) \\
(!A \lor !B)^N & \ident  \lnot\lnot (!A^N \lor !B^N) \\
(!A \lif !B)^N & \ident (!A^N \lif !B^N) \\
(\lforall[x][!A])^N & \ident \lforall[x][!A^N] \\
(\lexists[x][!A])^N & \ident \lnot\lnot\lexists[x][!A^N]
\end{align*}
Kolmogorov dan Glivenko mempunyai versi terjemahan ini untuk logika
proposisional; untuk logika predikat, terjemahan itu dikemukakan secara
terpisah oleh G\"odel dan Gentzen. Kita memperoleh

\begin{thm}
\begin{enumerate}
\item $!A \liff !A^N$ dapat dibuktikan secara klasik.
\item Jika $!A$ dapat dibuktikan secara klasik, maka $!A^N$ dapat dibuktikan
  secara intuisionistik.
\end{enumerate}
\end{thm}

Sekarang kita dapat membayangkan dialog berikut. Matematikawan klasik:
``Saya telah membuktikan $!A$!'' Matematikawan intuisionis: ``Bukti Anda tidak
valid. Yang sebenarnya Anda buktikan ialah $!A^N$.'' Matematikawan klasik:
``Bagi saya tidak masalah!'' Menurut matematikawan klasik, sang intuisionis
hanya memperdebatkan perbedaan remeh, karena keduanya ekuivalen. Namun, sang
intuisionis bersikeras bahwa terdapat perbedaan.

Perhatikan bahwa terjemahan di atas hanya menyangkut logika murni; terjemahan
itu tidak membahas pertanyaan mengenai aksioma \emph{nonlogis} yang tepat bagi
matematika klasik dan matematika intuisionistik, ataupun hubungan di antara
keduanya. Namun, perluasan kecil dari teorema di atas memberikan beberapa
informasi berguna:

\begin{thm}
Jika $\Gamma$ membuktikan $!A$ secara klasik, maka $\Gamma^N$ membuktikan
$!A^N$ secara intuisionistik.
\end{thm}

Dengan kata lain, jika $!A$ dapat dibuktikan secara klasik dari beberapa
hipotesis, maka $!A^N$ dapat dibuktikan dari terjemahan negasi ganda
hipotesis-hipotesis tersebut.

Untuk menunjukkan bahwa suatu kalimat atau !!{formula} proposisional valid
secara intuisionistik, kita hanya perlu memberikan bukti. Namun, bagaimana
kita dapat menunjukkan bahwa kalimat atau formula itu tidak valid? Untuk
tujuan tersebut, kita memerlukan semantik yang sahih dan sebaiknya lengkap.
Semantik yang dikemukakan oleh Kripke memenuhi kebutuhan itu dengan baik.

Kita dapat melakukan hal yang sama seperti bagi logika klasik: mendefinisikan
semantik, lalu membuktikan kesahihan dan kelengkapan. Namun, patut diperhatikan
perbedaan berikut. Dalam logika klasik, semantiknya merupakan semantik yang
``jelas'' dalam arti tertentu, tersirat dalam makna penghubung. Meskipun kita
dapat memberikan motivasi intuitif bagi semantik Kripke, semantik yang kedua
ini tidak memberikan kesan keniscayaan yang sama. Selain itu, gagasan mengenai
!!{structure} klasik merupakan gagasan matematis yang alami, sehingga kita
dapat memandang gagasan !!a{structure} sebagai perangkat untuk mengkaji logika
orde pertama klasik, atau memandang logika orde pertama klasik sebagai
perangkat untuk mengkaji !!{structure}s matematis. Sebaliknya, !!{structure}s
Kripke hanya dapat dipandang sebagai konstruksi logis; struktur itu tampaknya
tidak mempunyai kepentingan matematis yang mandiri.

Suatu !!{structure} Kripke~$\mModel M = \tuple{W, R, V}$ bagi bahasa
proposisional terdiri atas himpunan~$W$, urutan parsial~$R$ pada~$W$ yang
mempunyai !!{element} terkecil, dan penetapan ``monoton'' variabel proposisi
kepada !!{element}s dari~$W$. Intuisinya ialah bahwa !!{element}s dari $W$
mewakili ``dunia'', atau ``keadaan pengetahuan''; elemen $v \geq u$ mewakili
``keadaan masa depan yang mungkin'' dari~$u$; dan variabel proposisi yang
ditetapkan kepada~$u$ adalah proposisi yang diketahui benar dalam keadaan~$u$.
Relasi pemaksaan $\mSat{M}{!A}[w]$ kemudian memperluas hubungan ini ke
!!{formula}s sebarang dalam bahasa tersebut; bacalah $\mSat{M}{!A}[w]$
sebagai ``$!A$ benar dalam keadaan~$w$.'' Hubungan itu didefinisikan secara
induktif sebagai berikut:
\begin{enumerate}
\item $\mSat{M}{\Obj p_i}[w]$ jika dan hanya jika $\Obj p_i$ merupakan salah
  satu variabel proposisi yang ditetapkan kepada~$w$.
\item $\mSat/{M}{\lfalse}[w]$.
\item $\mSat{M}{(!A \land !B)}[w]$ jika dan hanya jika
  $\mSat{M}{!A}[w]$ dan $\mSat{M}{!B}[w]$.
\item $\mSat{M}{(!A \lor !B)}[w]$ jika dan hanya jika
  $\mSat{M}{!A}[w]$ atau $\mSat{M}{!B}[w]$.
\item $\mSat{M}{(!A \lif !B)}[w]$ jika dan hanya jika, setiap kali
  $w' \geq w$ dan $\mSat{M}{!A}[w']$, berlaku $\mSat{M}{!B}[w']$.
\end{enumerate}
Merupakan latihan yang baik untuk mencoba menunjukkan bahwa $\lnot (p \land
q) \lif (\lnot p \lor \lnot q)$ tidak valid secara intuisionistik, dengan
membangun suatu !!{structure} Kripke yang memberikan contoh penyangkal.

\end{document}
